\begin{definition}[\(\LJ\)-derivation]
An \(\LJ\)-derivation is a finite rooted tree of intuitionistic sequents.  A
leaf is the conclusion of an instance of \textsc{Id} or
\textsc{\(\bottom L\)}, and every nonleaf is the conclusion of exactly one
instance of one of the remaining displayed \(\LJ\) rules, whose premises are
its immediate predecessors.  A sequent is \(\LJ\)-derivable, written
\[
  \LJ\derives\Gamma\Rightarrow A,
\]
iff it is the root of such a tree.  The derivation is cut-free iff
none of its nodes is obtained by \textsc{Cut}.
\end{definition}

\begin{theorem}[Soundness of \(\LK\)]
Let \(\Gamma\) and \(\Delta\) be finite contexts.
If \(\LK\derives\Gamma\Rightarrow\Delta\), then every classical structure
and assignment satisfying every formula in \(\Gamma\) satisfies at least one
formula in \(\Delta\).
\end{theorem}

\paragraph{Proof sketch.}
Induct on the derivation.  Initial and structural cases follow directly from
the semantic reading of a sequent.  Each logical case unfolds the semantic
clause for its principal connective.  The quantifier cases use the substitution
lemma; the eigenvariable cases use the fact that the eigenvariable is absent
from the conclusion.  The individual case calculations are omitted.  Semantic
soundness of \(\LJ\) is stated only after its
intuitionistic semantics is defined in \cref{ch:intuitionistic-logic}.

